%%
%% Implementation Details (RGSW)
%% The specifics of the program that are not necessarily important for the math, and other (novel) optimizations
%%
\newpage\section{Implementation Details of the Core Library}
\label{sec:impl:core}
The core library, located in the \texttt{libs/core} directory of the git repository \cite{gitrepo}, extends OpenFHE~\cite{cryptoeprint:2022/915} with \gls{rgsw} functionality that the framework does not itself expose. % Its centerpiece is the abstract interface \texttt{IExtendedContext}, which augments an OpenFHE \texttt{CryptoContext} with the essential \gls{rgsw} operations: encryption, the external product, and the internal product.
This functionality is specified by an abstract interface, \texttt{IExtendedContext}, that any project can consume independently of the underlying implementation; of which the essential operations are listed in \autoref{tab:main-functions}. The default implementation realizes the interface through gadget decomposition of the RNS primes, as described in \autoref{sec:gadget-decomposition}.

%% TODO: Call it BV-impl instead of default impl
% Either here or in gadget decomposition somewhere: Following the convention of~\cite{cryptoeprint:2021/204, cryptoeprint:2022/915}, we refer to this family of decomposition-based techniques collectively as \emph{BV}. 

\begin{table}[H]
    \centering
    \caption{Primary functions declared by \texttt{IExtendedContext}}
    \label{tab:main-functions}
    \begin{tabular}{lcc}
        \toprule
        \textbf{Function Signature}           & \textbf{Returns} \\ % & \textbf{Definition}  \\
        \midrule
        \Call{EncryptRGSW}{pk, $\mu$}         & $\Call{RGSW}{\mu}$ \\ % & \ref{def:rgsw} \\
        \Call{EvalExternalProduct}{$x$, $C$}  & $x \boxdot C$    \\ % & \ref{def:prod:external} \\
        \Call{EvalInternalProduct}{$A$, $C$}  & $A \boxtimes C$  \\ % & \ref{def:prod:internal} \\
        \bottomrule
    \end{tabular}
\end{table}

Following standard software engineering practice, the interface fixes \emph{what} each operation computes but not \emph{how}, so code written against it runs unchanged under any conforming implementation.

%% TODO: Explain context -> Takes crypto context parameters that are considered global/known by all the algorithms without specifying them!
%% TODO: Implementation is hard-coupled to the BGV-RNS scheme. Maybe: As OpenFHE provides a unified interface for both BGV and BFV --- and the TFHE products are defined independent of scheme --- it would not require much work to add support for BFV to the core library.

\paragraph{Bitwise Operations}
The library takes $b = \log_2 B$ as the parameter and derives $B = 2^b$, so that the arithmetic of digit extraction reduces to constant-time bitwise operations (\autoref{tab:bitwise}); in particular, reduction modulo $B$ becomes a mask and division by $B$ a shift.

\begin{table}[H]
    \caption{Bitwise Operators}
    \label{tab:bitwise}
    \centering
    \begin{tabular}{lcc}
        \toprule
        \textbf{Named Function} & \textbf{Operator Notation} & \textbf{Value}   \\
        \midrule
        \Call{LeftShift}{$x, b$}  & $x \ll b$     & $x\cdot 2^b$ \\
        \Call{RightShift}{$x, b$} & $x \gg b$     & $\lfloor x/2^b \rfloor$ \\
        \Call{And}{$x, y$}        & $x \land y$   & $\sum_i x_i^{(2)}\cdot y_i^{(2)} \cdot 2^i$ \\
        \Call{Modulo}{$x, b$}     & $x \land [(1 \ll b) - 1]$ & $x \bmod{2^b}$ \\
        \bottomrule
    \end{tabular}
\end{table}

%%%
%%% Center residue
%%%
\begin{algorithm}[ht!]
\caption[Center]{\protect\Call{Center}{$x,b$}}
\label{alg:center}
\begin{algorithmic}[1]
    \Require $x, b \in \mathbb{Z}$
    \Ensure $[x]_{2^b} \in [-2^{b-1}, 2^{b-1})$
    \State mask $\gets \Call{LeftShift}{1, b} - 1$
    \State half $\gets \Call{LeftShift}{1, b-1}$
    \State\Return $\Call{And}{x + \text{half}, \text{mask}} - \text{half}$
\end{algorithmic}
\end{algorithm}

\newpage\paragraph{Bitwise Operations}
The library takes $b = \log_2 B$ as the parameter and derives $B = 2^b$, so that the arithmetic of digit extraction reduces to constant-time bitwise operations (\autoref{tab:bitwise}); in particular, reduction modulo $B$ becomes a mask and division by $B$ a shift.

\begin{table}[H]
    \caption{Bitwise Operators}
    \label{tab:bitwise}
    \centering
    \begin{tabular}{lcc}
        \toprule
        \textbf{Named Function} & \textbf{Operator Notation} & \textbf{Value}   \\
        \midrule
        \Call{LeftShift}{$x, b$}  & $x \ll b$     & $x\cdot 2^b$ \\
        \Call{RightShift}{$x, b$} & $x \gg b$     & $\lfloor x/2^b \rfloor$ \\
        \Call{And}{$x, y$}        & $x \land y$   & $\sum_i x_i^{(2)}\cdot y_i^{(2)} \cdot 2^i$ \\
        \Call{Modulo}{$x, b$}     & $x \land [(1 \ll b) - 1]$ & $x \bmod{2^b}$ \\
        \bottomrule
    \end{tabular}
\end{table}

%%%
%%% Center residue
%%%
\begin{algorithm}[ht!]
\caption[Center]{\protect\Call{Center}{$x,b$}}
\label{alg:center}
\begin{algorithmic}[1]
    \Require $x, b \in \mathbb{Z}$
    \Ensure $[x]_{2^b} \in [-2^{b-1}, 2^{b-1})$
    \State mask $\gets \Call{LeftShift}{1, b} - 1$
    \State half $\gets \Call{LeftShift}{1, b-1}$
    \State\Return $\Call{And}{x + \text{half}, \text{mask}} - \text{half}$
\end{algorithmic}
\end{algorithm}

\newpage%% TODO: Consider including automatic B selection if there is time
% The base $B = \lfloor \max q_i \rfloor$ is chosen automatically so to be the smallest $B$ that covers the largest RNS prime with $\ell$ digits when the program starts. This makes the gadget exact, even if the $q_i$ are generated at different orders (in OpenFHE the first prime is typically larger than the other primes for some reason).
% Takes only $\ell$ as an input parameter and computes $\log{B} = \left\lfloor \max\{\log(q_i)\}/\ell \right\rfloor + 1$

\subsection{RGSW Encryption}
\begin{itemize}
    \item \texttt{ExtendedContextBV} implements the double decomposition; digit decomposing each RNS limb
    \item Let $P_g(v) = g\otimes v$ and $P_\mathsf{BV}(v) = P_g \circ \hps{P}(v) = g\otimes (\hps{P}(v))$
\end{itemize}
\begin{gather*}
    C = Z + \mu G = Z + \mu \cdot (I_2\otimes  P_g(\hps{P}(1)) = Z + (I_2\otimes  P_g(\hps{P}(\mu)) \\
    P_g(\hps{P}(\mu)) = g\otimes ([\mu]_{q_0}, \dots, [\mu]_{q_{k-1}}) \\
    ([\mu g]_{q_0}, \dots [\mu g]_{q_{k-1}}) \text{ where } \mu g = (\mu, \mu B, \dots, \mu B^{\ell})\\
    ([\mu]_{q_0}, [\mu B]_{q_0}, \dots, [\mu B^\ell]_{q_{k-1}}) \in \mathbb{Z}_Q^{k\cdot \ell}
\end{gather*}

\begin{itemize}
    \item Observe that this is the same information as $(\mu, \mu B, \dots, \mu B^{\ell})$
    \item Because the $i$-th tower of $\mu$ is $[\mu]_{q_i}$, the HPS gadget row $j$ gives $[[\mu]_{q_i}]_{q_j} = 0 \forall i\not = j$
\end{itemize}

\begin{gather*}
    %([\mu]_{q_0}, [\mu B]_{q_0}, \dots, [\mu B^\ell]_{q_{k-1}}) \\
    \left[ \begin{array}{ccc|ccc|ccc}
        [\mu]_{q_0} & \cdots & [\mu B^{\ell -1 }]_{q_0} & 0 & \dots & 0 & 0 & \dots & 0 \\
        0 &  \dots & 0 &[\mu]_{q_1} & \cdots & [\mu B^{\ell -1 }]_{q_{1}} & 0 & \dots & 0 \\
        \vdots & \ddots & \vdots & \vdots & \ddots & \vdots & \vdots & \ddots & \vdots \\
        0 & \dots & 0 & 0 & \dots & 0 &[\mu]_{q_{k-1}} & \cdots & [\mu B^{\ell -1 }]_{q_{k-1}}  \\
    \end{array} \right]
\end{gather*}

\begin{itemize}
    \item Ignore all the 0 elements, then we get the exact definition of the RNS form of $\mu$, $\mu B$, and so on
    \item However, each row of the RGSW must still encode only 1 tower from each of the $\ell$ polynomials
\end{itemize}

\begin{equation*}
    {P}_{\mathsf{BV}}(\mu) \;\overset{\tiny\text{flatten}}{\longmapsto}\; 
    \begin{bmatrix}
        [\mu]_{q_0} & \cdots & [\mu B^{\ell - 1}]_{q_0} \\
        [\mu]_{q_1} & \cdots & [\mu B^{\ell - 1}]_{q_1} \\
        \vdots & \ddots & \vdots  \\
        [\mu]_{q_{k-1}} & \cdots & [\mu B^{\ell - 1}]_{q_{k-1}}  \\
    \end{bmatrix}
    = ([\mu]_{\mathcal{B}_Q}, \dots, [\mu B^{\ell - 1}]_{\mathcal{B}_Q})
\end{equation*}

\begin{itemize}
    \item Multiplying $\mu$ by a factor of $B^i$ for a single $i$ can be done in NTT
    \item Applying this to $Z + P_\mathsf{BV}(\mu)$ allows parallelization on $k\ell$ threads
    \item $\mathcal{B}_Q$, $\ell$ and $B$ are all known/global variables tied to the context
\end{itemize}

%%%
%%% Encrypt RGSW-BV
%%%
\begin{algorithm}
\caption{Parallel RGSW-BV Encryption}
\label{alg:bv:encrypt}
\begin{algorithmic}[1]
\Require Decomposition parameter $\ell$, base $B$
\Require RNS polynomial $[\mu]^{\mathsf{NTT}}_{\mathcal{B}_Q} \in \mathbb{Z}_{q_0} \times \dots \times  \mathbb{Z}_{q_{k-1}}$
\Ensure $\rgsw{(\mu)}$
\State Initialize empty array $C$ with $2k\ell$ slots
\For{$r \gets [0, k\ell)$ \textbf{in parallel}}
    \State $(i, j) \gets (r\bmod{\ell}, \lfloor r/\ell \rfloor)$
    \Comment{$i$-th power $B^i$, $j$-th tower}

    \State $mg \gets \Call{TowerMul}{\mu, j, B^i}$

    \State $C_r \gets \Call{Encrypt}{0} + (mg, 0) $
    \State $C_{r + k\ell} \gets \Call{Encrypt}{0} + (0, mg)$
    % FLAG: previous version indexed the second block as C_{r+\ell}; the offset is k\ell (matches the implementation).
\EndFor
\State \Return $C$
\end{algorithmic}
\end{algorithm}


\begin{itemize}
    \item Pre-computed value: $\Call{GetPower}{i, j} = B^i \bmod{q_j}$ where the modulo is just to keep numbers smaller in memory, negligible for performance 
    \item Then $\Call{TowerMul}{\mu, j, B^i}$ is really just extracting the $j-th$ limb, multiplying by the pre-computed value $B^i \bmod{q_j}$ calculate the product modulo $q_j$ before the tower is inserted back into $\mu$ 
\end{itemize}

\subsection{External Product}
%% TODO: Write about power-of-B gadget in earlier section!?
The decomposition on the power-of-$B$ gadget $g = (1, B, \dots, B^{\ell-1})$ is essentially digit extraction. The classical approach is to compute standard unsigned base-$B$ digits, producing a decomposition vector where each digit $d_i \in [0, B)$. Under our abstract gadget decomposition definition~\autoref{sec:gadget-decomposition}, the quality of this decomposition is $\beta = B - 1$.

From~\autoref{rem:extprod-asymmetry} we know the noise in the external product is a function of $\beta$. To exert tighter control over noise growth during external products, we naturally prefer to minimize this magnitude. By centering the digits into a signed representation where $d_i \in [-B/2, B/2)$, we halve the maximum absolute value, improving the overall decomposition quality to $\beta = B/2$. However, achieving this centered representation classically complicates parallelization. Standard signed recoding algorithms evaluate digits sequentially from least significant to most significant because shifting a positive digit $d_i \ge B/2$ into the negative range requires adding a $+1$ carry to the subsequent digit $d_{i+1}$. This ripple-carry dependency prevents us from slicing the integer into independent bit-windows, creating a computational bottleneck.

To avoid sequential carry propagation, we apply a well-known offset-binary trick by adding the constant $\varGamma$~\eqref{eq:carry-offset}.
\begin{equation}
    \label{eq:carry-offset}
    \varGamma = \frac{B}{2}\cdot\frac{B^{\ell} - 1}{B - 1}
\end{equation}

The second fraction is the standard mathematical identity for the sum of a finite geometric series, meaning~\eqref{eq:carry-offset} represents the $g$-decomposition of a base-$B$ number $\gamma$ where every digit is $B/2$.
\begin{gather*}
    \frac{B^{\ell} - 1}{B - 1} = \sum_{i=0}^{k-1} B^i = 1 + B + B^2 + \dots + B^{\ell - 1} \\
    \implies\;\; \varGamma = \sum_{i=0}^{k-1}\frac{B}{2}\cdot B^i = \langle \decomp{\gamma}{}, g \rangle
\end{gather*}

To decompose a positive integer $c \in \mathbb{Z}^{+}$ into signed digits $[d]_{B}$, we first apply the global offset to compute a shifted, non-negative integer $c' = c + \varGamma$. We then slice $c'$ into intermediate unsigned base-$B$ digits $u_i$. Because $B=2^b$ is a power of two, this extraction is performed purely via the independent bitwise \autoref{alg:bit:center}, completely bypassing sequential dependencies.

% \begin{algorithm}[ht!]
% \caption[Parallel Gadget Decomposition]{\protect\Call{ParallelDecompose}{$c, \ell, b$}}
% \label{alg:parallel-decomp}
% \begin{algorithmic}[1]
%     \Require Integer $c \in \mathbb{Z}$, decomposition length $\ell$, bit-length $b$
%     \Ensure Array of digits $d = (d_0, d_1, \dots, d_{\ell-1})$ where $d_i \in [-2^{b-1}, 2^{b-1})$ for $i < \ell - 1$
%     \State $\varGamma \gets \sum_{i=0}^{\ell-1} \Call{LeftShift}{1,\; b \cdot i + b - 1}$ \Comment{Precomputed constant offset}
%     \State $c' \gets c + \varGamma$ \Comment{Global shift}
%     \For{$i \gets 0$ \textbf{to} $\ell - 1$ \textbf{in parallel}}
%         \State $u_i \gets \Call{Modulo}{\Call{RightShift}{c',\; b \cdot i},\; b}$ \Comment{Extract independent bit-window}
%         \If{$i + 1 < \ell$}
%             \State $d_i \gets u_i - \Call{LeftShift}{1,\; b - 1}$ \Comment{Undo offset for lower digits}
%         \Else
%             \State $d_i \gets u_i$ \Comment{Top digit stays unsigned}
%         \EndIf
%     \EndFor
%     \State \Return $(d_0, d_1, \dots, d_{\ell-1})$
% \end{algorithmic}
% \end{algorithm}


% \begin{itemize}
%     \item \textbf{IMPORTANT:} decomposition determines the quality $\beta$, so the order of operations does matter for noise growth!
%     \item Decomposition used = per limb digit extraction
%     \item Must choose to decompose into $[-B/2, B/2)$ or $[0, B)$
%     \begin{itemize}
%         \item Centered has lower noise as the quality is $\beta = B/2$ but more complicated
%         \item Unsigned is far simpler and easily parallelizable but $\beta = B$
%     \end{itemize}
%     \item Centered representation creates an order; parallelization requires \textit{pre-processing}
% \end{itemize}


%%%
%%% Unsigned digit decomposition (optimized for B = 2^b)
%%% TODO: Move this to gadgets/BV section and have signed digit decomposition here
% \begin{algorithm}
% \caption{Decompose wrt. $B$ (\textcolor{red}{TODO: Signed digit decomposition})}
% \label{alg:bv:decompose}
% \begin{algorithmic}[1]
% \Require RNS polynomial $[x]_{\mathcal{B}_Q} \in \mathbb{Z}_{q_0} \times \mathbb{Z}_{q_1} \times \dots \times  \mathbb{Z}_{q_{k-1}}$, base $B = 2^b$, and decomposition parameter $\ell$
% \Ensure $\ell$ RNS ``digit'' polynomials $\mathbf{d} = (d_0, d_1, \dots, d_{k\ell-1}) \in \mathbb{Z}_B^{k\ell}$
% \State $a' \gets a$
% \For{$j = 0$ to $\ell - 1$}
%     \State $d_j \gets a' \pmod B$ \Comment{Extract the lower bits}
%     \If{$d_j \ge B/2$}
%         \State $d_j \gets d_j - B$ \Comment{Shift to signed range}
%     \EndIf
%     \State $a' \gets \frac{a' - d_j}{B}$ \Comment{Shift down and apply the carry}
% \EndFor
% \State \Return $\mathbf{d}$
% \end{algorithmic}
% \end{algorithm}

%%%
%%% 
%%%
% \begin{algorithm}
% \caption{Decompose wrt. $B$ (\textcolor{red}{TODO: Signed digit decomposition})}
% \label{alg:bv:decompose}
% \begin{algorithmic}[1]
% \Require Integer $[x]_{\mathcal{B}_Q} \in \mathbb{Z}_{q_0} \times \mathbb{Z}_{q_1} \times \dots \times  \mathbb{Z}_{q_{k-1}}$, base $B = 2^b$, and decomposition parameter $\ell$
% \Ensure $\ell$ RNS ``digit'' polynomials $\mathbf{d} = (d_0, d_1, \dots, d_{k\ell-1}) \in \mathbb{Z}_B^{k\ell}$
% \State $a' \gets a$
% \For{$j = 0$ to $\ell - 1$}
%     \State $d_j \gets a' \pmod B$ \Comment{Extract the lower bits}
%     \If{$d_j \ge B/2$}
%         \State $d_j \gets d_j - B$ \Comment{Shift to signed range}
%     \EndIf
%     \State $a' \gets \frac{a' - d_j}{B}$ \Comment{Shift down and apply the carry}
% \EndFor
% \State \Return $\mathbf{d}$
% \end{algorithmic}
% \end{algorithm}
\begin{algorithm}[ht!]
\caption[Parallel Signed Decomposition]{\protect\Call{Decompose}{$c, b, \ell$}}
\label{alg:decomp:parallel}
\begin{algorithmic}[1]
    \Require Integer $c \in \mathbb{Z}$, bit-length $b$ (where $B = 2^b$), decomposition parameter $\ell$
    \Ensure Array of signed digits $d = (d_0, d_1, \dots, d_{\ell-1})$ with quality $\beta = 2^{b-1}$
    
    \State $\varGamma \gets \sum_{i=0}^{\ell-1} 2^{b-1} \cdot 2^{b \cdot i}$ \Comment{This offset is typically precomputed}
    \State $c' \gets c + \varGamma$ \Comment{Apply the global shift}
    \State mask $\gets \Call{LeftShift}{1, b} - 1$
    \State half $\gets \Call{LeftShift}{1, b-1}$
    
    \For{$i \gets 0$ \textbf{to} $\ell - 1$ \textbf{in parallel}}
        \State $u_i \gets \Call{And}{\Call{RightShift}{c',\; i \cdot b},\; \text{mask}}$ \Comment{Extract the intermediate unsigned digit}
        \If{$i + 1 < \ell$}
            \State $d_i \gets u_i - \text{half}$ \Comment{Undo the offset to center the digit}
        \Else
            \State $d_i \gets u_i$ \Comment{The most significant digit remains unsigned}
        \EndIf
    \EndFor
    
    \State\Return $d = (d_0, d_1, \dots, d_{\ell-1})$
\end{algorithmic}
\end{algorithm}

Once we have a 

\begin{itemize}
    \item Uses the standard external product function \autoref{alg:external-product}
    \item Or a slightly parallelized version: the decompose function runs in parallel, not much gain from parallizing further, the additions are already cheap/fast-ish.
\end{itemize}

%%%
%%% External Product
%%%
\begin{algorithm}
\caption{External Product $\rlwe \boxdot \rgsw$ (Naive)}
\label{alg:external-product}
\begin{algorithmic}[1]
\Require RLWE ciphertext $(c_0, c_1) \in R_Q^2$, RGSW ciphertext $C \in R_Q^{2\ell}$
\Ensure RLWE ciphertext of product
\State $acc \gets (0, 0)$
\State $d_0 \gets \Call{Decompose}{c_0}$
\State $d_1 \gets \Call{Decompose}{c_1}$
\For{$r \gets [0, \ell)$}
    \State $acc[0] \gets acc[0] + d_0[r] \cdot C[r][0] + d_1[r]\cdot C[r+\ell][0]$
    \State $acc[1] \gets acc[1] + d_0[r] \cdot C[r][1] + d_1[r]\cdot C[r+\ell][1]$
\EndFor
\State \Return $acc$
\end{algorithmic}
\end{algorithm}

However, our implementation contains a series of optimizations aimed at . It can be divided into three distinct stages, each stage depending on the previous stage...

%%%
%%% Implemented (staged, parallel) external product
%%%
\begin{algorithm}
\caption{Parallel External Product $\rlwe \boxdot \rgsw$ (RNS)}
\label{alg:bv:extprod}
\begin{algorithmic}[1]
\Require Decomposition parameter $\ell$, base $B = 2^b$, offset $\delta = \frac{B}{2}\cdot\frac{B^{\ell-1}-1}{B-1}$
\Require RLWE ciphertext $(c_0, c_1)$, each $[c_m]^{\mathsf{NTT}}_{\mathcal{B}_Q} \in \mathbb{Z}_{q_0} \times \dots \times  \mathbb{Z}_{q_{k-1}}$
\Require RGSW ciphertext $C$ with $2k\ell$ rows as produced by \autoref{alg:bv:encrypt}
\Ensure RLWE ciphertext of the product

\For{$(m, j) \gets [0, 2) \times [0, k)$ \textbf{in parallel}}
    \Comment{Stage 1: $2k$ single-tower iNTTs}
    \State $a_{m,j} \gets \Call{iNTT}{c_{m,j}}$
    \Comment{tower $j$ of $c_m$, coefficient form}
\EndFor

\For{$(s, t) \gets [0, 2k\ell) \times [0, k)$ \textbf{in parallel}}
    \Comment{Stage 2: slice, lift, $2k^2\ell$ single-tower NTTs}
    \State $(m, r) \gets (\lfloor s/k\ell \rfloor,\ s \bmod k\ell)$
    \Comment{digit slot $s$ pairs with row $s$ of $C$}
    \State $(i, j) \gets (r \bmod \ell,\ \lfloor r/\ell \rfloor)$
    \State $u \gets \lfloor (a_{m,j} + \delta) / B^i \rfloor \bmod B$
    \Comment{coefficient-wise shift-and-mask}
    \If{$i + 1 < \ell$}
        \State $u \gets u - B/2$
        \Comment{undo offset; the top digit stays unsigned}
    \EndIf
    \State $d_{s,t} \gets \Call{NTT}{[u]_{q_t}}$
    \Comment{signed lift into $\mathbb{Z}_{q_t}$}
\EndFor

\For{$(m, t) \gets [0, 2) \times [0, k)$ \textbf{in parallel}}
    \Comment{Stage 3: $2k$ independent accumulators}
    \State $acc_{m,t} \gets \sum_{s=0}^{2k\ell - 1} d_{s,t} \odot C[s][m]_t$
    \Comment{fused multiply-accumulate in $\mathbb{Z}_{q_t}$}
\EndFor
\State \Return $(acc_0, acc_1)$
\end{algorithmic}
\end{algorithm}


\subsection{Internal Product}
The internal product is trivially implemented as repeated external products as per the definition~\autoref{def:prod:internal}.

%% TODO: Include algorithmic literal line-by-line as impl

\subsection{SIMD Batching}
Nothing in the gadget or in either product depends on how the plaintext ring is interpreted, so the implementation inherits the slot packing of \autoref{def:simd} from OpenFHE unchanged. A caller that builds its plaintexts with \texttt{MakePackedPlaintext} rather than \texttt{MakeCoefPackedPlaintext} obtains \gls{rgsw} ciphertexts whose message is a vector of $N$ independent slots, and every external and internal product then acts on all $N$ slots at once. The library selects the matching zero plaintext automatically, so no other change is required.

% The cost of a product is unchanged by this. A packed ciphertext is the same object of the same size, and the $2k\ell$ rows are decomposed and accumulated exactly as before, which is what makes the packing worthwhile. An application holding $N$ independent instances of the same binary computation can place one in each slot and evaluate all of them for the price of one, reducing the amortized cost per instance by a factor of $N$. The gain is in throughput and never in latency, and it is available only to applications whose work decomposes into independent instances in this way.